\documentclass[11pt,a4paper]{article}

% --- Encoding & language ---------------------------------------------------
\usepackage[utf8]{inputenc}
\usepackage[T1]{fontenc}
\usepackage[french]{babel}
\usepackage{lmodern}
\usepackage{microtype}

% --- Layout ----------------------------------------------------------------
\usepackage[a4paper, top=2.4cm, bottom=2.4cm, left=2.5cm, right=2.5cm]{geometry}
\usepackage{setspace}
\onehalfspacing

% --- Visuals ---------------------------------------------------------------
\usepackage{xcolor}
\definecolor{kvink}{RGB}{12,18,28}
\definecolor{kvgold}{RGB}{196,142,42}
\definecolor{kvalert}{RGB}{226,58,58}
\definecolor{kvgreen}{RGB}{63,163,108}
\definecolor{kvgrey}{RGB}{120,128,140}
\definecolor{kvbg}{RGB}{246,246,242}

\usepackage{titlesec}
\titleformat{\section}
  {\color{kvink}\Large\bfseries\sffamily}
  {\thesection.}{0.6em}{}
\titleformat{\subsection}
  {\color{kvink}\large\bfseries\sffamily}
  {\thesubsection}{0.6em}{}
\titleformat{\subsubsection}
  {\color{kvgold}\normalsize\bfseries\sffamily}
  {}{0pt}{}

\usepackage{booktabs}
\usepackage{array}
\usepackage{tabularx}
\renewcommand{\arraystretch}{1.25}

\usepackage{enumitem}
\setlist{leftmargin=1.4em, itemsep=2pt, topsep=3pt}

\usepackage{fancyhdr}
\pagestyle{fancy}
\fancyhf{}
\fancyhead[L]{\sffamily\color{kvink}\textbf{KOVER.IA}}
\fancyhead[R]{\sffamily\color{kvgrey}\small Hackathon AI Paris~2026}
\fancyfoot[C]{\sffamily\color{kvgrey}\small \thepage}
\renewcommand{\headrulewidth}{0.4pt}

\usepackage{tcolorbox}
\tcbuselibrary{skins,breakable}
\newtcolorbox{insight}[1][]{
  enhanced, breakable, colback=kvbg, colframe=kvink,
  boxrule=0.5pt, arc=2pt, left=10pt, right=10pt, top=8pt, bottom=8pt,
  fontupper=\sffamily\small, #1
}
\newtcolorbox{quote_kv}[1][]{
  enhanced, colback=white, colframe=kvgold,
  boxrule=1pt, arc=0pt, left=14pt, right=14pt, top=10pt, bottom=10pt,
  borderline west={3pt}{0pt}{kvgold},
  fontupper=\itshape, #1
}

\usepackage[hidelinks]{hyperref}

% ===========================================================================
\begin{document}

\begin{titlepage}
\centering
\vspace*{2.5cm}
{\sffamily\Huge\bfseries\color{kvink} KOVER.IA}\\[0.2cm]
{\sffamily\Large\color{kvgold} \rule{4cm}{1pt}}\\[0.4cm]
{\sffamily\LARGE\bfseries Rapport projet complet}\\[0.6cm]
{\sffamily\large\color{kvgrey} Problématique, solution, architecture,\\[0.1cm]
positionnement et stratégie de commercialisation}\\[3.2cm]

\begin{tcolorbox}[
  enhanced, colback=kvbg, colframe=kvink,
  boxrule=0.5pt, arc=2pt, width=12cm,
  left=12pt, right=12pt, top=10pt, bottom=10pt,
  fontupper=\sffamily
]
\centering
\textbf{Hackathon AI Paris 2026}\\[0.4cm]
\small Détection en temps réel du blanchiment crypto\\
et des attaques flashloan par intelligence artificielle.
\end{tcolorbox}

\vfill
{\sffamily\color{kvgrey}\small Avril 2026 \quad\textbar\quad Version 1.0 \quad\textbar\quad \texttt{kover.ia}}
\end{titlepage}

\renewcommand{\contentsname}{Table des matières}
{\sffamily\tableofcontents}
\thispagestyle{fancy}
\newpage

% ===========================================================================
\section{Problématique : quel problème règle-t-on ?}

\subsection{Un secteur DeFi structurellement vulnérable}

La finance décentralisée (DeFi) gère aujourd'hui plus de
\textbf{120~milliards de dollars} de valeur verrouillée à travers des
protocoles ouverts opérant 24~heures sur~24 sur des blockchains
publiques. En 2024, selon Chainalysis, plus de
\textbf{2,2~milliards de dollars} ont été volés sur des protocoles
DeFi --- une augmentation de 21~\% par rapport à 2023. À ces vols
directs s'ajoute un phénomène encore plus coûteux pour l'écosystème~:
le \textbf{blanchiment des fonds dérobés}, qui se déroule en quelques
minutes via des mixeurs, des bridges cross-chain et des plateformes
d'échange centralisées peu régulées.

\subsection{Le problème central : la vélocité}

Le blanchiment crypto suit un patron stéréotypé observé sur tous les
hacks majeurs récents (Euler Finance, Ronin Bridge, Wormhole,
Curve Finance)~:

\begin{enumerate}
  \item \textbf{Splitting} : éclatement sur 4 à 12 wallets en moins de
        deux minutes.
  \item \textbf{Mixing} : 5 à 20 transactions intermédiaires via
        Tornado Cash, Railgun ou des wallets jetables.
  \item \textbf{Bridging} : transit cross-chain (BSC, Polygon, Arbitrum)
        via Wormhole, Stargate, Synapse.
  \item \textbf{Off-ramp} : dépôt sur un CEX peu régulé pour conversion
        en monnaie fiat.
\end{enumerate}

L'ensemble de cette séquence se déroule en
\textbf{15 à 90 minutes}. Aucune équipe humaine, même expérimentée, ne
peut suivre ce rythme. Lorsqu'un analyste ouvre l'investigation, les
fonds sont déjà dans un mixeur ou sur un bridge, et la traçabilité
on-chain devient exponentiellement plus coûteuse.

\subsection{Une seconde menace : les flashloan attacks}

En parallèle du blanchiment post-vol, les protocoles DeFi sont la cible
d'\textbf{attaques par flashloan}~: emprunt non collatéralisé d'une
somme massive, manipulation d'un oracle ou d'un pool de liquidité, et
remboursement dans la même transaction. Tout se joue dans le mempool
public, en quelques centaines de millisecondes avant le minage du bloc.
Là encore, aucun outil grand public ne couvre la fenêtre d'interception.

\begin{quote_kv}
Le problème central que nous adressons est donc double~: un problème de
\textbf{détection} (voir le blanchiment ou l'attaque émerger en temps
réel) et un problème de \textbf{décision} (produire une recommandation
actionnable assez vite pour qu'elle soit utile).
\end{quote_kv}

% ===========================================================================
\section{Limites des solutions actuelles sur le marché}

Les solutions de surveillance AML on-chain existantes
(Chainalysis Reactor, TRM Labs, Elliptic Lens, Forta) sont
historiquement bâties sur trois piliers~:

\begin{itemize}
  \item \textbf{Listes d'adresses étiquetées}, maintenues manuellement
        par des équipes d'analystes (sanctions OFAC, mixeurs connus,
        wallets de hackers).
  \item \textbf{Heuristiques de propagation} fondées sur des règles
        métier déterministes (nombre de \emph{hops}, montants,
        fréquence).
  \item \textbf{Tableaux de bord d'investigation} optimisés pour
        l'analyse rétrospective, sur des fenêtres de jours ou semaines.
\end{itemize}

\subsection{Quatre limites structurelles}

\subsubsection{1. Une approche rétrospective, pas prédictive}

Ces outils \emph{étiquettent} après les faits. Ils ne modélisent pas
le graphe de transactions comme un système dynamique et ne prédisent
jamais la prochaine étape probable d'un flux suspect.

\subsubsection{2. Aucune représentation latente apprise}

Les wallets sont définis par des \emph{labels} manuels. Un wallet
inconnu mais structurellement suspect (proche topologiquement de
mixeurs) ne sera jamais signalé tant qu'un analyste ne l'aura pas
classé.

\subsubsection{3. Une latence en minutes, pas en secondes}

Les rapports Chainalysis se mesurent en heures, le scoring TRM en
minutes. Cette latence rend la \textbf{réponse opérationnelle}
impossible~: les fonds sont déjà loin.

\subsubsection{4. Black-box non auditable}

Les décisions algorithmiques sont opaques. Avec MiCA et l'AMLR-2024,
les régulateurs européens exigent désormais une traçabilité
algorithmique vérifiable que les acteurs en place ne fournissent qu'au
prix d'audits externes coûteux.

\begin{insight}
\textbf{Le vide à combler.} Il n'existe \emph{aucun} produit grand
public combinant représentations apprises sur graphe, prédiction
séquentielle de la trajectoire des fonds, narratif LLM exploitable,
et provenance cryptographique vérifiable --- en moins de cinq secondes.
\end{insight}

% ===========================================================================
\section{Notre solution : KOVER.IA}

\subsection{Philosophie produit}

\textsc{KOVER.IA} est un \textbf{capteur prédictif temps réel} qui
observe le flux de transactions Ethereum, modélise le graphe émergent,
prédit la trajectoire des fonds suspects, et produit un rapport
d'incident en moins de cinq secondes. La solution combine deux modules
synergiques~:

\begin{itemize}
  \item \textbf{Module AML} : détection et prédiction du blanchiment
        sur le graphe des transactions confirmées.
  \item \textbf{Module BFD} (\emph{Behavioral Flow Detection}) :
        sentinel mempool spécialisé dans l'interception des flashloan
        attacks avant minage.
\end{itemize}

\subsection{Ce que nous délivrons concrètement}

\begin{enumerate}
  \item Un score de fraude par wallet, calculé toutes les 5~secondes
        sur un graphe roulant.
  \item Une prédiction de la prochaine destination probable
        (Uniswap, Binance, Hyperliquid\ldots).
  \item Un rapport d'incident en langage naturel, streamé
        token-par-token, contenant un résumé exécutif, une analyse
        technique et des actions recommandées.
  \item Un \emph{halt} automatique en cas de flashloan critique
        détectée dans le mempool.
  \item Un manifest cryptographique signé HMAC-SHA256 prouvant la
        chaîne de provenance de chaque inférence.
\end{enumerate}

% ===========================================================================
\section{Architecture technique}

\subsection{Pipeline multi-modèle orchestré}

\begin{tabularx}{\linewidth}{@{}lXl@{}}
\toprule
\textbf{Modèle} & \textbf{Rôle} & \textbf{Fréquence} \\
\midrule
GAT (Graph Attention Network) &
Score de fraude par wallet ; capte les dépendances structurelles
du graphe (un wallet propre individuellement peut être condamné par
ses voisins). &
Toutes les 5~s \\
\addlinespace
LSTM séquentiel &
Prédit la prochaine destination du flux à partir des cinq derniers
wallets du chemin de propagation pondéré par le score GAT. &
Toutes les 6~s \\
\addlinespace
Qwen~3 235B (Cerebras) &
Génère un rapport d'incident structuré en streaming token-par-token. &
Toutes les 30~s \\
\bottomrule
\end{tabularx}

\medskip

Le pipeline ingère un flux Ethereum-like calibré sur le dataset
\emph{Salam Ammari} (73~034 wallets, 71~250 transactions, ratio
fraude $\approx$~3,7~\%). Un graphe roulant de 50 nœuds est maintenu
en mémoire sur quatre tiers (source du hack, splitters, mixeurs,
terminaux) --- topologie fidèle aux hacks observés.

\subsection{Module Behavioral Flow Detection (BFD)}

Le sentinel \texttt{kover-bfd-mev}~:

\begin{itemize}
  \item Pré-filtre des sélecteurs de méthodes (Aave, dYdX, Balancer,
        Maker).
  \item Simulation \texttt{eth\_call} pour évaluer la transaction sans
        diffusion.
  \item Verdict IA Cerebras (Qwen) classant la sévérité
        (\emph{LOW / MEDIUM / CRITICAL}) et l'\emph{exploit class}.
  \item Diffusion via Server-Sent Events vers un dashboard temps réel.
\end{itemize}

Performance mesurée~: \textbf{6,5~M événements/s}, latence bout-en-bout
de l'ordre de la \textbf{dizaine de millisecondes}.

\subsection{Provenance cryptographique}

Chaque inférence est attachée à un manifest signé HMAC-SHA256 contenant
l'empreinte SHA-256 des fichiers de poids, les métadonnées
d'entraînement, la liste des inférences récentes avec latences. Un
auditeur externe peut recalculer le HMAC et vérifier qu'aucune
prédiction n'a été falsifiée.

\subsection{Métriques de performance}

\begin{tabularx}{\linewidth}{@{}lXl@{}}
\toprule
\textbf{Métrique} & \textbf{Détail} & \textbf{Valeur} \\
\midrule
GAT --- accuracy validation & 14k wallets test & \textbf{97,0~\%} \\
GAT --- recall fraude & classe minoritaire & \textbf{94,0~\%} \\
LSTM --- accuracy par classe & clean / Scam / Phish & 89 / 91 / 58~\% \\
Cerebras --- débit & Qwen~3 235B streaming & 49~tokens/s \\
Pipeline complet & cycle bout-en-bout & < 5~s \\
Couverture tests & 80 tests pytest & 95--100~\% \\
\bottomrule
\end{tabularx}

% ===========================================================================
\section{Concurrence et positionnement}

\subsection{Cartographie du marché}

\begin{tabularx}{\linewidth}{@{}lXl@{}}
\toprule
\textbf{Acteur} & \textbf{Positionnement} & \textbf{Cible} \\
\midrule
Chainalysis & Compliance \& forensique & Banques, gouvernements \\
TRM Labs & Risk scoring quasi temps réel & Exchanges, fintechs \\
Elliptic & AML transaction monitoring & Institutions financières \\
Forta Network & Détection on-chain décentralisée & Protocoles, DAO \\
Hexagate / Hypernative & Détection d'attaques DeFi & Équipes sécurité protocoles \\
\bottomrule
\end{tabularx}

\subsection{Comparatif fonctionnel}

\begin{tabularx}{\linewidth}{@{}lcccc@{}}
\toprule
& \textbf{Chainalysis} & \textbf{TRM} & \textbf{Forta} & \textbf{KOVER.IA} \\
\midrule
Score IA appris (GNN)            & \texttimes & \texttimes & partiel & \checkmark \\
Prédiction destination (LSTM)    & \texttimes & \texttimes & \texttimes & \checkmark \\
Rapport LLM streamé              & \texttimes & \texttimes & \texttimes & \checkmark \\
Manifest crypto signé            & \texttimes & \texttimes & \texttimes & \checkmark \\
Sentinel flashloan               & \texttimes & partiel    & \checkmark & \checkmark \\
Latence médiane d'alerte         & minutes    & secondes   & secondes   & < 5~s \\
Open source partiel              & \texttimes & \texttimes & partiel    & \checkmark \\
\bottomrule
\end{tabularx}

\subsection{Pourquoi nous : quatre axes différenciants}

\subsubsection{1. La triple couche IA (graphe + séquence + langage)}

Aucun concurrent ne combine aujourd'hui un GNN, un RNN séquentiel et un
LLM en un pipeline temps réel orchestré. Chaque modèle existe
isolément~; leur \emph{orchestration} est notre apport propriétaire.

\subsubsection{2. Provenance cryptographique native}

Le manifest signé HMAC répond directement à l'exigence MiCA / AMLR-2024
de traçabilité algorithmique. Là où les concurrents répondent par audit
externe, nous répondons par preuve vérifiable.

\subsubsection{3. Latence sub-5~secondes bout-en-bout}

Cette latence déplace le cas d'usage de l'\emph{investigation
post-mortem} vers la \emph{réponse opérationnelle}. C'est un changement
de catégorie produit, pas une amélioration incrémentale.

\subsubsection{4. Surface d'intégration ouverte}

WebSocket public + endpoint REST signé. N'importe quelle DAO, oracle ou
multi-sig peut consommer nos verdicts sans contrat commercial préalable
--- modèle natif crypto, contrairement aux SDK propriétaires des acteurs
établis.

% ===========================================================================
\section{Coûts de production}

Le \emph{COGS} d'une instance \textsc{Defend} en production est dominé
par trois postes~:

\begin{itemize}
  \item \textbf{Inférence Cerebras} --- $\approx$ 7~\$/M tokens. Sur un
        pattern 1 rapport/30~s, 250 tokens/rapport, 24$\times$7 :
        $\approx$ 720~M tokens/an, soit \textbf{5~040~\$/an}.
  \item \textbf{Hébergement backend} --- AWS m6a.large + Redis,
        $\approx$ 80~\$/mois en prod, 3 zones DR :
        \textbf{2~880~\$/an}.
  \item \textbf{Données on-chain} --- QuickNode \emph{Build}
        (mempool + archive) à 99~\$/mois, soit
        \textbf{1~188~\$/an} par chaîne.
\end{itemize}

\medskip

\textbf{COGS consolidé : $\approx$ 9~100~\$/an} pour un client
\textsc{Defend} qui paie 144~k\$/an. Soit une \textbf{marge brute
> 93~\%} --- typique d'un SaaS B2B deeptech, et qui finance la R\&D ML,
l'expansion produit (Solana, Bitcoin) et le sales.

% ===========================================================================
\section{Comment on vend : à qui, pourquoi}

\subsection{À qui : nos quatre segments cibles}

\begin{tabularx}{\linewidth}{@{}lXl@{}}
\toprule
\textbf{Segment} & \textbf{Besoin principal} & \textbf{ARR cible} \\
\midrule
Protocoles DeFi (TVL > 100~M\$) & Surveillance flashloan + réponse auto & 60--250~k\$ \\
Exchanges régulés (CEX) & Filtre AML temps réel sur dépôts entrants & 120--500~k\$ \\
DAO de gouvernance & Alerte indépendante pour vote on-chain & 24--80~k\$ \\
Régulateurs / TRACFIN & Capacité d'audit avec preuve crypto & 250~k\$+ \\
\bottomrule
\end{tabularx}

\subsection{Pourquoi : la fenêtre d'opportunité}

\begin{enumerate}
  \item \textbf{Réglementaire} --- MiCA entre pleinement en vigueur fin
        2026. Les protocoles DeFi opérant en Europe \emph{devront}
        prouver leur conformité AML.
  \item \textbf{Économique} --- une attaque flashloan moyenne coûte
        25~M\$ au protocole victime. Notre offre \textsc{Guard} à
        30~k\$/an représente 0,12~\% du coût d'un seul incident
        évité.
  \item \textbf{Technologique} --- la disponibilité de l'inférence
        Cerebras (49 tok/s sur Qwen 235B) rend possible aujourd'hui
        une orchestration impensable il y a 18~mois.
\end{enumerate}

\subsection{Stratégie commerciale : land \& expand}

Trois plans SaaS, du gratuit au sur-devis~:

\begin{tabularx}{\linewidth}{@{}lXll@{}}
\toprule
\textbf{Plan} & \textbf{Inclus} & \textbf{Tarif} & \textbf{Engagement} \\
\midrule
\textsc{Watch} &
WebSocket public, scoring GAT, 50 alertes/jour. & 0 \$/mois & Aucun \\
\addlinespace
\textsc{Guard} &
Pipeline complet, 5~000 alertes/jour, sentinel 1 chaîne, SLA 24~h. &
2~500 \$/mois & 12 mois \\
\addlinespace
\textsc{Defend} &
Quotas illimités, 5 chaînes, SLA 2~h, manifest persisté, fine-tuning. &
12~000 \$/mois & 12--36 mois \\
\addlinespace
\textsc{Enterprise} &
On-premise / VPC, audit OFAC, compliance MiCA, formations. &
sur devis & 24--36 mois \\
\bottomrule
\end{tabularx}

\medskip

Le segment \emph{protocoles DeFi} est notre \textbf{point d'entrée
prioritaire}~: sensibilité au risque maximale, budget sécurité existant
(audits Trail of Bits, primes Immunefi), cycle de décision court (5--15
personnes décisionnaires).

\subsection{Go-to-market en trois temps}

\begin{enumerate}
  \item \textbf{Design partners} (mois 0--12) : 3 à 5 protocoles à
        60~k\$/an, validation produit, témoignages.
  \item \textbf{Inbound DAO + Forta marketplace} (mois 6--18) :
        plan \textsc{Watch} gratuit pour génération de leads.
  \item \textbf{Sales enterprise} (mois 12+) : équipe AE/SE pour
        cycles CEX et régulateurs.
\end{enumerate}

% ===========================================================================
\section{Pourquoi nous : crédibilité du projet livré}

\textsc{KOVER.IA} n'est pas un slide-deck~: le projet livré au
Hackathon AI Paris~2026 est \textbf{pleinement fonctionnel}.

\begin{itemize}
  \item \textbf{Modèles entraînés} sur données Ethereum réelles (dataset
        Salam Ammari), avec accuracy validation 97~\% sur le GAT.
  \item \textbf{Pipeline orchestré} GAT + LSTM + LLM Cerebras, latence
        bout-en-bout < 5~s.
  \item \textbf{Sentinel flashloan} opérationnel sur le mempool
        Ethereum, débit 6,5~M~ev/s.
  \item \textbf{Manifest cryptographique} exposé publiquement,
        vérifiable HMAC-SHA256.
  \item \textbf{Qualité industrielle} : GitHub Actions, SonarCloud,
        80 tests pytest, couverture 95--100~\% sur les modules
        critiques.
\end{itemize}

\begin{quote_kv}
La fenêtre d'opportunité est étroite. MiCA s'applique fin 2026, et les
protocoles DeFi qui n'auront pas adopté un AML traçable et auditable
seront mécaniquement exclus des marchés européens. \textsc{KOVER.IA}
se positionne comme l'option \emph{native crypto} pour répondre à
cette contrainte, là où Chainalysis et Elliptic restent ancrés dans un
modèle de compliance bancaire hérité.
\end{quote_kv}

\begin{insight}
\textbf{Synthèse.} Nous résolvons un problème réel (le blanchiment et
les flashloans en temps réel), avec une approche techniquement inédite
(triple couche IA + provenance crypto), pour une cible solvable
(protocoles DeFi, CEX, DAO, régulateurs), sur une fenêtre
réglementaire ouverte (MiCA), avec une marge SaaS structurelle (>~93~\%)
et un produit déjà fonctionnel et auditable.
\end{insight}

\end{document}
